\section{Spacetime}

Spacetime is not in the base. It is what the base becomes at scale, and each piece of it was measured against a control that could have failed. \cite{pollard2026vibetest}

The base has a finite symmetry, the group of $1152$, with no continuous rotations. Coarse-grain the tones and that finite symmetry swells toward the continuous $\mathrm{SO}(8)$, the way a fine enough polygon rounds into a circle, and with it come the rotations and boosts of relativity. The test is sharp. On the twenty-four directions the dispersion is isotropic, the same in every direction, to fourth order, while the flat cubic lattice shows a preferred frame and fails. Lorentz invariance is emergent and selective, and the hyperbolic substrate passes where the flat one does not.

\begin{result}[Emergent Lorentz invariance]
The discrete $F_4$ symmetry restores to continuous $\mathrm{SO}(8)$ in the infrared, and the $24$-direction dispersion is isotropic to fourth order. A flat cubic lattice breaks isotropy and carries a preferred frame. \cite{pollard2026vibetest}
\end{result}

The speed limit is exact and early. A tone moves one dock per beat, so nothing outruns one dock per tick, and that finite light cone is present before any physics is read off. Coarse-grain the motion and the massless dispersion is straight, $\omega = c\,k$, matching the continuum.

The discreteness leaves a faint, testable mark. Because the symmetry is only restored at scale, the smallest distances carry a tiny violation of it, growing as the square of the energy divided by the Planck scale, far too small to notice at ordinary energies but in principle visible in the most energetic light that crosses the universe. The bound from gamma-ray bursts already rules out a flat lattice, whose violation grows at first order in the energy, while the hyperbolic substrate, violating only at second order, passes comfortably. It is a real astrophysical test the model survives and a naive discretization fails.

The residual has a precise origin worth naming. The bulk's own twenty-four directions are isotropic to fourth order, the good geometry, but matter does not propagate in the bulk, it propagates on the cubic cusp, and a cubic lattice is anisotropic at fourth order, so the small leftover preferred-frame effect is the cusp's and is not shielded by the bulk's symmetry. It is the one genuinely testable fingerprint of the discreteness, around eight percent between an axis and a diagonal at the lattice scale and falling as the square of the wavelength toward longer waves, which places it near one part in $10^{40}$ at laboratory energies, comfortably inside the current bounds. That our own universe is itself only approximately flat, to a fraction of a percent, and finite in age, makes an exactly-flat-only-in-the-limit discrete space a match rather than a contradiction.

Time has a definite shape, a growing block. The past is the settled interior, fixed and done. The future is open at the wake, not yet laid down. The present is the growing edge where new docks are born at peace. There is no block of all-time sitting complete, only a crystal still forming, and the boundary where it forms carries no reversibility, which is where the arrow lives. Space and time are not the same kind of thing at the base. Space is the mesh, four directions of it. Time is the growth of the mesh. They fuse only at scale, as spacetime, which is why the two look so different up close and so unified far away.
